\documentclass[12pt,a4paper]{article}
\usepackage[T2A]{fontenc}
\usepackage[utf8]{inputenc}
\usepackage[russian]{babel}
\usepackage{amsmath,amssymb,amsthm}
\usepackage{geometry}
\geometry{margin=2.5cm}
\usepackage{graphicx}
\usepackage{hyperref}
\hypersetup{colorlinks=true,linkcolor=blue,citecolor=blue,urlcolor=blue}
\usepackage{booktabs}
\usepackage{enumitem}

\newtheorem{theorem}{Теорема}
\newtheorem{remark}{Замечание}

\title{\textbf{Единая геометродинамическая модель элементарных частиц\\
как дефектов упругого эфира:\\
дислокационные кольца, плёнка Ферми и четыре взаимодействия}}
\author{А.~Ю.~Дроздов}
\date{24 октября 2026 г.}

\begin{document}
\maketitle

\begin{abstract}
\noindent
Представлено последовательное изложение геометродинамической модели, в которой
элементарные частицы трактуются как топологические дефекты (дислокационные и
дисклинационные кольца) упругого эфира, а фундаментальные взаимодействия --- как
различные проекции единого упругого отклика среды. Модель строится на калибровочной
теории дислокаций Кадича--Эделена и Лазара \cite{lazar2010,kadic1983}, где магнитное
поле отождествляется с плотностью дислокаций, а сила Пича--Кёлера играет роль силы
Лоренца. Гравитация вводится через механизм <<плёнки Ферми>> --- деформируемой
границы интеграла действия, натянутой ортогонально искривлённым мировым трубкам
ускоряющихся зарядов (обобщение работы Ферми 1923~г.\ о парадоксе $4/3$
\cite{fermi1923}). Показано, что закон обратных квадратов для гравитации является
теоремой (следствием функции Грина лапласиана), а не подгонкой. Модель формулируется
в упругом базисе единиц $(L,T,M)$, без привлечения электрического заряда как
первичной сущности. Честно очерчены границы: выводятся \emph{структура} взаимодействий
(законы спадания, знаки, иерархия) и ряд соотношений между масштабами; абсолютные
значения фундаментальных констант ($\alpha$, $G$) остаются входными параметрами.
Работа носит характер рабочего препринта и предназначена для первого знакомства с темой.

\medskip
\noindent\textbf{Ключевые слова:} калибровочная теория дислокаций, эфир, плёнка Ферми,
магнитный монополь, вектор Бюргерса, сила Пича--Кёлера, парадокс $4/3$, геометродинамика.
\end{abstract}

\tableofcontents
\newpage

\section{Введение}

\subsection{Мотивация: частица как дефект среды}
Идея о том, что элементарные частицы могут быть не точечными объектами, а
протяжёнными структурами в некоторой среде, имеет давнюю историю --- от вихрей
Кельвина до современных солитонных моделей. Настоящая работа развивает
конкретную версию этой программы: \emph{частица есть топологический дефект
упругого континуума} (<<эфира>>), а именно замкнутая дислокационная линия
(кольцо) или её обобщения. Такой подход привлекателен тем, что опирается на
хорошо разработанный математический аппарат теории дефектов в кристаллах
\cite{hirth1982,kadic1983} и его калибровочную формулировку \cite{lazar2010}.

\subsection{Почему дислокации}
В кристаллах дислокации --- это линейные дефекты, характеризуемые вектором
Бюргерса $\mathbf{b}$, который измеряет невязку обхода вокруг линии дефекта.
М.~Лазар \cite{lazar2010} показал, что динамическая калибровочная теория
трансляционной группы $T(3)$ приводит к уравнениям, структурно тождественным
уравнениям Максвелла (таблица~\ref{tab:corr}). Это даёт физически осмысленный
<<словарь>> между теорией дефектов и электродинамикой.

\begin{table}[h]
\centering
\caption{Соответствие Максвелл~$\leftrightarrow$~теория дислокаций (по Лазару \cite{lazar2010}).}
\label{tab:corr}
\begin{tabular}{ll}
\toprule
Электродинамика & Теория дислокаций \\
\midrule
Магнитное поле $\mathbf{B}$ & плотность дислокаций $T^a$ \\
Электрическое поле $\mathbf{E}$ & ток дислокаций $I^a$ \\
Магнитное возбуждение $\mathbf{H}$ & псевдомоментное напряжение $H_a$ \\
Магнитный поток $\oint\mathbf{B}\,d\mathbf{S}$ & вектор Бюргерса $\oint T^a = b$ \\
Сила Лоренца & сила Пича--Кёлера \\
\bottomrule
\end{tabular}
\end{table}

\subsection{Ключевая особенность: заряд магнитный, не электрический}
В предлагаемой модели \emph{электрический заряд не является первичной сущностью}.
Первичен магнитный заряд, отождествляемый с потоком плотности дислокаций, то
есть с вектором Бюргерса. Электрон строится как прецессирующий магнитный диполь.
Это соответствует давней линии, связывающей дислокации и дираковские монополи
\cite{mistura1995}.

\subsection{Структура работы}
В разделе~\ref{sec:units} вводится упругий базис единиц. Раздел~\ref{sec:particle}
описывает частицу как прецессирующий диполь. Раздел~\ref{sec:fermi} излагает
механизм плёнки Ферми и вывод гравитации. Раздел~\ref{sec:channels} сводит шесть
каналов взаимодействия. Раздел~\ref{sec:limits} честно очерчивает границы модели.

\section{Упругий базис единиц}
\label{sec:units}
Мы строим систему единиц из трёх первооснов упругой среды: длины $L$, времени $T$
и массы $M$. Электрического заряда и ампера в базисе \emph{нет}. Основные величины:
\begin{itemize}[nosep]
\item вектор Бюргерса $b$ (смещение решётки) --- размерность $[L]$;
\item модуль сдвига $\mu$ --- размерность $[M L^{-1}T^{-2}]$ (Па);
\item плотность $\rho$ --- $[M L^{-3}]$;
\item скорость звука/света $c=\sqrt{\mu/\rho}$ --- $[L T^{-1}]$.
\end{itemize}
Магнитный заряд как поток плотности дислокаций $q_m=\oint T^a\,dS$ имеет
размерность длины $[L]$ --- это внутренне согласованно в упругой теории
(в отличие от СИ, где магнитный заряд измеряется в А$\cdot$м; перенос единиц СИ
на упругую среду не требуется и приводит к путанице).

\section{Частица как прецессирующий диполь}
\label{sec:particle}

\subsection{Геометрия}
Электрон моделируется как пара дислокационных колец (радиуса $a$), плоскости
которых \emph{неподвижны}, а вдоль каждого кольца с частотой $\omega$ бежит
<<особая точка>> --- локализованный магнитный заряд $q_m=\Delta b$ (разрыв
вектора Бюргерса). Под <<прецессией>> понимается именно движение особой точки по
неподвижному кольцу, а не качание плоскостей. Два кольца разнесены на расстояние
$d$ и несут противоположные заряды $\pm q_m$ (N- и S-полюса), образуя магнитный
диполь.

\subsection{Масса из самодействия}
Масса не вводится отдельно, а есть энергия упругого самодействия дефекта:
\begin{equation}
m = \frac{\mu\, b^2 a}{c^2},
\end{equation}
что имеет правильную размерность $[M]$.

\subsection{Магнитный момент и спин}
Магнитный момент есть \emph{статический} дипольный момент $\mu_{\text{маг}}=q_m\,d$
(а не момент кругового тока: ток по кольцу магнитный, и наивная формула
токовой петли здесь неприменима). Спин связан с топологией: пара обычных колец
имеет период конфигурации $2\pi$ (бозон), тогда как пара мёбиусовых колец
(кольцо + анти-кольцо) --- период $4\pi$, что отвечает фермиону со спином $1/2$.
В мёбиус-паре два кольца несут \emph{противоположные} направления вектора
Бюргерса ($+\Delta b$ и $-\Delta b$, полюса $N$ и $S$) --- отсюда статический
магнитный диполь. Существенно, что твист Бюргерса $b(\varphi)=\cos(\varphi/2)$
есть поворот Франка на $\pi$ за обход, то есть \emph{дисклинационная} ($SO(3)$)
степень свободы: мёбиусов сектор модели уже лежит в $SO(3)$, а магнитный момент
$\mu=\nu=n/2$ выводится как поток спин-связности (Осипов). Гидродинамически тот
же объект --- вихрь: завихренность отвечает плотности Франка, а обязательная
для устойчивости прецессия ($\omega>0$) и есть циркуляция этого вихря.

\subsection{Устойчивость и аннигиляция}
Противоположные заряды $\pm q_m$ притягиваются и без движения схлопнулись бы.
Прецессия создаёт центробежный барьер (сохранение момента импульса
$J=m a^2\omega$), стабилизирующий пару на конечном расстоянии:
\begin{equation}
E(d) = -\frac{A}{d} + \frac{J^2}{2 m d^2},\qquad d_{\text{eq}}=\frac{J^2}{mA}.
\end{equation}
При $\omega\to0$ барьер исчезает --- пара аннигилирует (рекомбинация Френкеля,
суммарный вектор Бюргерса обращается в нуль).

\section{Плёнка Ферми и вывод гравитации}
\label{sec:fermi}

\subsection{Исторический контекст: парадокс $4/3$}
Э.~Ферми в 1923~г.\ \cite{fermi1923} показал, что известное расхождение между
электродинамической ($\tfrac{4}{3}U/c^2$) и релятивистской ($U/c^2$) массами
заряженного тела устраняется, если границу интегрирования действия проводить
\emph{строго ортогонально} искривлённой мировой трубке ускоряющегося заряда.
Эту деформируемую границу мы называем \emph{плёнкой Ферми}. Плёнка ---
не слагаемое лагранжиана, а граница области интегрирования действия. Поэтому
она не является динамическим полем, не запаздывает и мгновенно подстраивается
под искривление мировых трубок. Запаздывает и излучает только само поле,
получаемое из минимизации действия. Плёнку Ферми следует отличать от граничного члена Гиббонса--Хокинга--Йорка в
общей теории относительности \cite{hehl1995}. Наша плёнка отличается от него тем,
что не входит в лагранжиан: член Гиббонса--Хокинга--Йорка есть слагаемое
действия и подчиняется динамике поля, а плёнка Ферми --- лишь граница области
интегрирования. Поэтому она не запаздывает и мгновенно подстраивается под
искривления мировых трубок.

\subsection{Источник деформации}
Ускоряющийся заряд искривляет свою мировую трубку; плёнка, оставаясь
ортогональной, деформируется. Плотность источника деформации пропорциональна
\emph{квадрату} собственного ускорения $W=\omega^2 a$:
\begin{equation}
\sigma = \mu\, b^2\,\frac{\langle W^2\rangle}{c^2} > 0 .
\end{equation}
Знак положителен всегда (квадрат), поэтому гравитация --- всегда притяжение и
не зависит от знака заряда (универсальность).

\subsection{Закон обратных квадратов как теорема}
\begin{theorem}[Закон $1/R^2$]
Если деформация плёнки $\Delta$ подчиняется уравнению равновесия упругой
мембраны с точечным источником,
\begin{equation}
\gamma\,\nabla^2\Delta(\mathbf r) = -\sigma\,\delta^3(\mathbf r-\mathbf r_0),
\end{equation}
то решением является функция Грина лапласиана, $\Delta(R)=\sigma/(4\pi\gamma R)$,
откуда потенциал взаимодействия $U(R)\propto 1/R$ и сила $F(R)\propto 1/R^2$.
\end{theorem}
Показатель степени следует из трёхмерности пространства, а не из подгонки.

\subsection{Параметры жёсткости: два масштаба}
\paragraph{Что такое $\mu_{\text{эфир}}$.}
Величина $\mu_{\text{эфир}}$ --- это \emph{модуль сдвига вакуума}, то есть
жёсткость эфира по отношению к сдвиговой деформации, с размерностью давления
$[\text{Па}]$. Именно через этот упругий модуль в базисе $(L,T,M)$ (без
си-ампера) определяются основные величины модели:
\begin{align}
c &= \sqrt{\mu_{\text{эфир}}/\rho} && (\text{скорость поперечной волны}),\\
m &= \mu_{\text{эфир}}\,b^2 a / c^2 && (\text{масса частицы}),\\
F &= \mu_{\text{эфир}}\,b^2 / R \cdot (\text{геометрия}) && (\text{кулоновская сила}).
\end{align}
Физически $\mu_{\text{эфир}}$ отвечает за сдвиговые деформации (дислокации), из
которых в модели построены заряд и масса; поэтому он задаёт физику на масштабе
частицы (порядка комптоновской длины $\lambda_C$). Численное значение
получено \emph{калибровкой}: подстановкой массы электрона $m_e$ и его
классического радиуса $r_e$ в формулу самоэнергии дислокационного кольца и
разрешением относительно $\mu_{\text{эфир}}$. Это \emph{не} вывод из первооснов,
а фиксация одного числа по электрону, что отмечается явно. Более поздний разбор массы
(инерция сосредоточена в особых точках, их энергия задаётся модулем $c_1$ Лазара)
показал, что это значение --- лишь оценка порядка; см.\ расширенную версию.

Жёсткость плёнки имеет размерность давления $[\text{Па}]$. Проверка на модели
показывает, что \emph{один} модуль не может отвечать и за массу/кулон, и за
гравитацию: требуемые значения различаются на $\sim\!84$ порядка. Поэтому
необходимы два масштаба:
\begin{align}
\mu_{\text{эфир}} &\approx 2{,}7\times10^{28}~\text{Па} && (\text{ЭМ/масса, комптоновский масштаб}),\\
\gamma_P &= \frac{c^7}{8\pi\hbar G^2}\approx 1{,}8\times10^{112}~\text{Па} && (\text{гравитация, планковский масштаб}).
\end{align}
Гравитационная постоянная $G$ входит в $\gamma_P$: она является \emph{входным}
параметром жёсткости эфира, а не выводится из более фундаментальных величин.

\paragraph{Что такое $\gamma_P$.}
Величина $\gamma_P$ --- это \emph{планковская жёсткость вакуума}, второй модуль
упругости эфира, отвечающий за гравитационный отклик (деформацию плёнки Ферми),
с той же размерностью давления $[\text{Па}]$, что и $\mu_{\text{эфир}}$. Её
комбинация
\begin{equation}
\gamma_P = \frac{c^7}{8\pi\hbar G^2}
\end{equation}
построена так, чтобы жёсткость имела размерность давления и воспроизводила
гравитацию на планковском масштабе: в неё входят скорость света $c$, квант
действия $\hbar$ и гравитационная постоянная $G$. Подчеркнём честно: $G$ здесь
--- \emph{входной} параметр, а не выводимый; $\gamma_P$ не даёт $G$ из более
фундаментальных величин, а наоборот, вбирает её в себя. Физически $\gamma_P$
отвечает за тот же упругий отклик вакуума, но в гравитационном секторе, где
характерный масштаб --- планковская длина $\ell_P=\sqrt{\hbar G/c^3}$, а не
комптоновская. Колоссальное отношение
$\gamma_P/\mu_{\text{эфир}}\sim10^{84}$ есть прямое выражение слабости
гравитации по сравнению с кулоновским/массовым сектором и составляет главную
нерешённую проблему модели (происхождение иерархии двух масштабов).

\section{Шесть каналов взаимодействия}
\label{sec:channels}
Все взаимодействия суть проекции единого упругого отклика. Их законы спадания и
источники сведены в таблицу~\ref{tab:channels}; степенные законы --- структурные
следствия (функция Грина, диполь-дипольное взаимодействие, масса переносчика,
конфайнмент), а не результат подгонки.

\begin{table}[h]
\centering
\caption{Шесть каналов взаимодействия в модели.}
\label{tab:channels}
\begin{tabular}{llll}
\toprule
Канал & Закон & Источник & Константа связи \\
\midrule
Кулон & $1/R^2$ & произведение зарядов-Бюргерсов & $\alpha=7{,}30\times10^{-3}$ \\
Магнитное & $1/R^4$ & диполь--диполь $q_m d$ & $\sim\alpha/4$ \\
Гравитация & $1/R^2$ & $\langle W^2\rangle$ плёнки Ферми & $\alpha_G=1{,}75\times10^{-45}$ \\
Обмен & $1/R^3$ & фаза прецессии $\cos2\Delta\phi$ & $\sim\alpha$ \\
Слабое & $e^{-R/R_W}$ & киральность твиста & $\sim 1/30$ \\
Сильное & $\propto R$ & частичный $b/3$ + лента & $\sim 1$ \\
\bottomrule
\end{tabular}
\end{table}

Вся система замыкается на два безразмерных числа:
\begin{equation}
\alpha = \frac{k e^2}{\hbar c}=\frac{1}{137{,}04},\qquad
\alpha_G = \left(\frac{m_e}{m_{\text{Planck}}}\right)^2 = 1{,}75\times10^{-45},
\end{equation}
откуда слабость гравитации относительно кулона: $\alpha_G/\alpha\approx 2{,}4\times10^{-43}$.

\section{Границы модели: что доказано, что принято, что открыто}
\label{sec:limits}
В духе научной честности мы явно разделяем уровни достоверности.

\paragraph{Доказано (структурно).} Законы спадания всех каналов; притяжение и
универсальность гравитации; размерная замкнутость массы; согласованность
магнитного заряда как величины размерности длины в упругом базисе; механизм
аннигиляции и защита прецессией.

\paragraph{Принято (входные параметры).} Константы $\alpha$ и $\alpha_G$;
гравитационная постоянная $G$ (вложена в жёсткость $\gamma_P$); наличие двух
масштабов жёсткости и их отношение.

\paragraph{Открыто.} Происхождение численных значений $\alpha$ и $\alpha_G$;
строгий вывод гиромагнитного отношения $g=2$ (в классической картине колец
получается $g=1$; $g=2$ требует спинорного уровня описания в духе
\cite{osipov2010}); размерный множитель в источнике гравитации, требующий
динамической (зависящей от времени) жёсткости плёнки; численные значения
констант слабого и сильного секторов.

\section{Заключение}
Предложена внутренне согласованная модель, в которой частицы суть упругие
дефекты эфира, а взаимодействия --- проекции единого отклика среды через силу
Пича--Кёлера и деформацию плёнки Ферми. Модель объясняет \emph{структуру}
(законы, знаки, иерархию взаимодействий), не претендуя на вывод фундаментальных
констант из первых принципов. Такая постановка делает теорию фальсифицируемой:
например, для магнитно-гравитационных устройств она предсказывает ничтожную
($\sim\!10^{-13}$) относительную поправку к весу, что несовместимо с заявлениями
о макроскопической антигравитации. Дальнейшая работа --- строгий спинорный вывод
$g=2$, уточнение динамической жёсткости плёнки и перевод на английский для
архивной публикации.

\begin{thebibliography}{99}
\bibitem{fermi1923} E.~Fermi, \emph{Correzione di una contraddizione tra la teoria
elettrodinamica e quella relativistica delle masse elettromagnetiche},
Nuovo Cimento \textbf{25}, 159 (1923).
\bibitem{lazar2010} M.~Lazar, \emph{On the fundamentals of the three-dimensional
translation gauge theory of dislocations}, Math.\ Mech.\ Solids \textbf{16}(3), 253 (2010).
\bibitem{kadic1983} A.~Kadi\'c, D.~G.~B.~Edelen, \emph{A Gauge Theory of Dislocations
and Disclinations}, Lecture Notes in Physics \textbf{174}, Springer (1983).
\bibitem{lazar2009} M.~Lazar, \emph{On the Higgs mechanism and stress functions in the
translational gauge theory of dislocations}, arXiv:0903.0990 (2009).
\bibitem{lazar2009b} M.~Lazar, C.~Anastassiadis, \emph{The gauge theory of dislocations:
static solutions of screw and edge dislocations}, arXiv:0802.0670 (2009).
\bibitem{osipov2010} E.~A.~Kochetov, V.~A.~Osipov, \emph{Dirac fermions on a disclinated
flexible surface}, arXiv:1001.4450 (2010).
\bibitem{osipov1999} E.~A.~Kochetov, V.~A.~Osipov, \emph{Gauge theory of disclinations on
fluctuating elastic surfaces}, arXiv:math-ph/9903018 (1999).
\bibitem{hehl1995} F.~W.~Hehl, J.~D.~McCrea, E.~W.~Mielke, Y.~Ne'eman, \emph{Metric-affine
gauge theory of gravity}, Phys.\ Rep.\ \textbf{258}, 1 (1995); arXiv:gr-qc/9402012.
\bibitem{hirth1982} J.~P.~Hirth, J.~Lothe, \emph{Theory of Dislocations}, 2nd ed.,
Wiley (1982).
\bibitem{mistura1995} L.~Mistura, \emph{Crystal dislocations and Dirac monopoles},
Int.\ J.\ Eng.\ Sci.\ \textbf{33}, 2149 (1995).
\bibitem{rohrlich1960} F.~Rohrlich, \emph{Self-energy and stability of the classical
electron}, Am.\ J.\ Phys.\ \textbf{28}, 639 (1960).
\end{thebibliography}

\end{document}
